% =====================================================================
% Appendix A: Source Code and Supplementary Material
% =====================================================================
\chapter{Source Code and Supplementary Material}
\label{ch:appendix}

The complete implementation of both pipelines developed for this
thesis, covering the image-based and video-based systems described in
Chapters~\ref{ch:methodology} and~\ref{ch:system-design}, has been made
publicly available as open-source code; the repository link is provided
in the Abstract. The repository contains the following materials:

\begin{itemize}[noitemsep]
    \item \textbf{\texttt{jersey\_recognition\_images.ipynb}} --- Complete
    image pipeline notebook: SoccerNet data loading, YOLOv8x detection,
    YOLOv8x-Pose torso localization, multi-variant OCR pipeline with
    dynamic zoom and shape-based digit disambiguation, ResNet-18
    legibility classifier training, and player-level evaluation.

    \item \textbf{\texttt{jersey\_recognition\_videos.ipynb}} --- Complete
    video pipeline notebook: YOLOv5 football detection, ByteTrack
    multi-object tracking, temporal role voter for referee
    identification, PARSeq (with EasyOCR fallback) digit recognition on
    torso crops, temporal jersey voting (60-frame window), and
    ground-truth evaluation against MOT and jersey-number CSV files.

    \item \textbf{\texttt{output/}} --- All output figures, CSV manifests,
    and evaluation reports produced by both pipelines.

    \item \textbf{\texttt{gt/}} --- Manually annotated ground-truth
    jersey-number CSV files for all ten evaluated broadcast clips.
\end{itemize}

Key implementation modules and their roles are summarized in
Table~\ref{tab:code-modules}, and representative source-code excerpts
for each are given in Sections~\ref{sec:appx-image-code}
and~\ref{sec:appx-video-code}. Excerpts are lightly trimmed (error
handling and logging statements omitted) for print length; the complete,
runnable versions are in the notebooks listed above.

\begin{table}[H]
    \centering
    \caption{Key implementation modules referenced in this thesis.}
    \label{tab:code-modules}
    \small
    \begin{tabular}{@{}p{4.5cm}p{3cm}p{5.5cm}@{}}
        \toprule
        \textbf{Module / Class} & \textbf{Pipeline} & \textbf{Role} \\
        \midrule
        \texttt{clean\_main\_image()} & Image & Denoising, CLAHE contrast, and unsharp masking \\
        \texttt{PlayerDetector} & Image & YOLOv8x person detection with fallback cascade \\
        \texttt{PoseHelper} & Image & YOLOv8x-Pose keypoint extraction and torso derivation \\
        \texttt{ocr\_digits\_strict()} & Image & Multi-variant OCR: zoom, binarize, split, shape vote \\
        \texttt{LegibilityModel} & Image & ResNet-18 binary legibility classifier \\
        \texttt{JerseyRecognizer} & Video & PARSeq recognizer with automatic EasyOCR fallback \\
        \texttt{BYTETracker} & Video & Standalone ByteTrack multi-object tracker \\
        \texttt{TemporalRoleVoter} & Video & Referee identification via semantic class history \\
        \texttt{TemporalJerseyVoter} & Video & 60-frame sliding-window jersey number aggregation \\
        \texttt{GroundTruthTrackingAndJerseyEvaluator} & Video & Ground-truth tracking and OCR accuracy evaluator \\
        \bottomrule
    \end{tabular}
\end{table}

\section{Image Pipeline --- Selected Code Excerpts}
\label{sec:appx-image-code}

\subsection{Global Image Cleaning}

\begin{lstlisting}[caption={\texttt{clean\_main\_image()} --- denoising, CLAHE contrast, and unsharp masking applied to every input image before detection (Section~\ref{sec:preprocessing}).}]
def clean_main_image(img_bgr):
    """Denoise + CLAHE contrast + edge-preserving smoothing + unsharp mask."""
    if img_bgr is None or img_bgr.size == 0:
        return img_bgr
    dn = cv2.fastNlMeansDenoisingColored(img_bgr, None, 5, 5, 7, 15)
    lab = cv2.cvtColor(dn, cv2.COLOR_BGR2LAB)
    L, A, B = cv2.split(lab)
    L = cv2.createCLAHE(clipLimit=1.6, tileGridSize=(8, 8)).apply(L)
    cont = cv2.cvtColor(cv2.merge([L, A, B]), cv2.COLOR_LAB2BGR)
    smooth = cv2.bilateralFilter(cont, d=0, sigmaColor=18, sigmaSpace=6)
    blur = cv2.GaussianBlur(smooth, (0, 0), 1.0)
    sharp = cv2.addWeighted(smooth, 1.35, blur, -0.35, 0)
    return sharp
\end{lstlisting}

\subsection{Player Detection and Pose Estimation}

\begin{lstlisting}[caption={\texttt{PlayerDetector} and \texttt{PoseHelper} --- YOLOv8x person detection and YOLOv8x-Pose keypoint extraction, each with a cascading fallback across model sizes (Section~\ref{sec:models}).}]
class PlayerDetector:
    def __init__(self, weights=None):
        self.model = None
        tried = ([weights] if weights else []) + \
                ["yolov8x.pt", "yolov8l.pt", "yolov8m.pt", "yolov8s.pt", "yolov8n.pt"]
        for w in tried:
            try:
                self.model = YOLO(w)
                print(f"[det] Using {w}")
                break
            except Exception as e:
                print(f"[det] Failed {w}: {e}")

    def detect_players(self, image_path, conf=PLAYER_CONF):
        if self.model is None:
            return np.empty((0, 4))
        r = self.model(image_path, conf=conf, verbose=False, max_det=50)[0]
        b = r.boxes.xyxy.detach().cpu().numpy() \
            if r.boxes is not None and r.boxes.xyxy is not None else np.empty((0, 4))
        if r.boxes is not None and r.boxes.cls is not None:
            cls = r.boxes.cls.detach().cpu().numpy().astype(int)
            b = b[cls == 0]  # class 0 = person
        return b if b.size else np.empty((0, 4))

class PoseHelper:
    def __init__(self, weights=None):
        self.model = None
        tried = ([weights] if weights else []) + \
                ["yolov8x-pose.pt", "yolov8l-pose.pt", "yolov8m-pose.pt",
                 "yolov8s-pose.pt", "yolov8n-pose.pt"]
        for w in tried:
            try:
                self.model = YOLO(w)
                print(f"[pose] Using {w}")
                break
            except Exception as e:
                print(f"[pose] Failed {w}: {e}")

    def infer(self, image_path, conf=POSE_CONF):
        if self.model is None:
            return {"boxes": np.empty((0, 4)), "kpts": None}
        r = self.model(image_path, conf=conf, verbose=False, max_det=50)[0]
        b = r.boxes.xyxy.detach().cpu().numpy() \
            if r.boxes is not None and r.boxes.xyxy is not None else np.empty((0, 4))
        k = None
        if hasattr(r, "keypoints") and r.keypoints is not None \
                and r.keypoints.xy is not None:
            k = r.keypoints.xy.detach().cpu().numpy()
        return {"boxes": b if b.size else np.empty((0, 4)), "kpts": k}
\end{lstlisting}

\subsection{Legibility Classifier and Training Configuration}

\begin{lstlisting}[caption={\texttt{LegibilityModel} --- ResNet-18 binary legibility classifier with class-imbalance-weighted BCE loss and the \texttt{ReduceLROnPlateau} scheduler configuration (factor 0.6) used in Section~\ref{sec:legibility-training}.}]
class LegibilityModel(nn.Module):
    def __init__(self, pretrained=True, dropout=0.25):
        super().__init__()
        weights = torchvision.models.ResNet18_Weights.DEFAULT if pretrained else None
        backbone = torchvision.models.resnet18(weights=weights)
        num_f = backbone.fc.in_features
        backbone.fc = nn.Sequential(nn.Dropout(dropout), nn.Linear(num_f, 1))
        self.backbone = backbone

    def forward(self, x):
        return self.backbone(x).squeeze(1)  # raw logits

model = LegibilityModel(pretrained=True).to(DEVICE)

# pos_weight = (#negative / #positive) to counter class imbalance
num_pos = sum(train_labels_l)
num_neg = len(train_labels_l) - num_pos
pos_weight = torch.tensor([num_neg / max(1, num_pos)], dtype=torch.float32).to(DEVICE) \
    if num_pos > 0 else torch.tensor([1.0], dtype=torch.float32).to(DEVICE)

criterion = nn.BCEWithLogitsLoss(pos_weight=pos_weight)
optimizer = torch.optim.AdamW(model.parameters(), lr=2e-4, weight_decay=1e-4)
scheduler = torch.optim.lr_scheduler.ReduceLROnPlateau(
    optimizer, mode="min", factor=0.6, patience=2
)
\end{lstlisting}

\section{Video Pipeline --- Selected Code Excerpts}
\label{sec:appx-video-code}

\subsection{PARSeq Recognizer with Automatic EasyOCR Fallback}

\begin{lstlisting}[caption={\texttt{JerseyRecognizer} --- attempts to load PARSeq at start-up and falls back to EasyOCR if unavailable (Section~\ref{sec:models}). This is the exact fallback mechanism illustrated in Figure~\ref{fig:video-step6-ocr}.}]
class JerseyRecognizer:
    def __init__(self, device="cpu"):
        self.device = device
        self.backend = None
        self.model = None
        self.img_transform = None

        try:
            self.model = torch.hub.load(
                "baudm/parseq", "parseq",
                pretrained=True, trust_repo=True,
            ).eval().to(device)

            from strhub.data.module import SceneTextDataModule
            self.img_transform = SceneTextDataModule.get_transform(
                self.model.hparams.img_size
            )
            self.backend = "parseq"
            print("[OCR] Using PARSeq backend.")

        except Exception as e:
            print(f"[OCR] PARSeq unavailable ({e}); falling back to EasyOCR.")

        if self.backend is None:
            import easyocr
            self.model = easyocr.Reader(["en"], gpu=(device != "cpu"))
            self.backend = "easyocr"
            print("[OCR] Using EasyOCR backend.")

    @staticmethod
    def _clean(text):
        digits = re.sub(r"\D", "", str(text))
        if not digits:
            return None
        digits = digits[:2]
        try:
            v = int(digits)
        except Exception:
            return None
        return str(v) if 0 <= v <= 99 else None

    def recognize(self, roi_bgr):
        if roi_bgr is None or roi_bgr.size == 0:
            return None, 0.0

        if self.backend == "parseq":
            from PIL import Image
            rgb = cv2.cvtColor(roi_bgr, cv2.COLOR_BGR2RGB)
            pil = Image.fromarray(rgb)
            with torch.no_grad():
                inp = self.img_transform(pil).unsqueeze(0).to(self.device)
                logits = self.model(inp)
                probs = logits.softmax(-1)
                preds, confs = self.model.tokenizer.decode(probs)
            text = preds[0] if preds else ""
            conf = (float(np.mean(confs[0].cpu().numpy()))
                    if len(confs) and len(confs[0]) else 0.0)
            num = self._clean(text)
            return (num, conf) if num else (None, 0.0)

        # EasyOCR fallback path
        res = self.model.readtext(
            roi_bgr, allowlist="0123456789", detail=1, paragraph=False,
        )
        best, best_conf = None, 0.0
        for _, text, conf in res:
            n = self._clean(text)
            if n and conf > best_conf:
                best, best_conf = n, float(conf)
        return (best, best_conf) if best else (None, 0.0)
\end{lstlisting}

\subsection{Temporal Role and Jersey Voting}

\begin{lstlisting}[caption={\texttt{TemporalRoleVoter} and \texttt{TemporalJerseyVoter} --- confidence-weighted temporal aggregation for referee exclusion and jersey-number consensus (Section~\ref{sec:temporal-voting}, Equation~\ref{eq:temporal-vote}).}]
class TemporalRoleVoter:
    def __init__(self, window=20, ref_ratio=0.6, min_votes=5):
        self.window = window
        self.ref_ratio = ref_ratio
        self.min_votes = min_votes
        self.history = defaultdict(lambda: deque(maxlen=window))
        self._locked = set()

    def observe(self, tid, class_name):
        if tid is None:
            return
        tid = int(tid)
        self.history[tid].append(1 if class_name == "referee" else 0)
        obs = self.history[tid]
        if len(obs) >= self.min_votes:
            r = sum(obs) / len(obs)
            if r >= self.ref_ratio:
                self._locked.add(tid)
            elif tid in self._locked and r < self.ref_ratio * 0.5:
                self._locked.discard(tid)

    def is_referee(self, tid):
        return tid is not None and int(tid) in self._locked


class TemporalJerseyVoter:
    def __init__(self, window=60, min_conf=0.25, min_votes=3):
        self.window = window
        self.min_conf = min_conf
        self.min_votes = min_votes
        self.history = defaultdict(lambda: deque(maxlen=window))
        self._final = {}

    def add(self, tid, number, conf):
        if tid is None or number is None or conf < self.min_conf:
            return
        self.history[int(tid)].append((str(number), float(conf)))

    def vote(self, tid):
        if tid is None:
            return None
        tid = int(tid)
        obs = self.history.get(tid)
        if not obs or len(obs) < self.min_votes:
            return self._final.get(tid)
        w = Counter()
        for number, conf in obs:
            w[str(number)] += float(conf)
        best = w.most_common(1)[0][0]
        self._final[tid] = best
        return best
\end{lstlisting}

\subsection{Ground-Truth Tracking and Jersey-Number Evaluator}

\begin{lstlisting}[caption={\texttt{GroundTruthTrackingAndJerseyEvaluator} --- initialization and accumulator structure for the ground-truth evaluation described in Section~\ref{sec:gt-eval}.}]
class GroundTruthTrackingAndJerseyEvaluator:
    def __init__(self, csv_path, iou_threshold=0.30, frame_offset=0):
        self.csv_path = csv_path
        self.iou_threshold = float(iou_threshold)
        self.frame_offset = int(frame_offset)
        self.available = False

        self.gt_by_frame = defaultdict(list)
        self.gt_total_boxes = 0

        # OCR evaluation accumulators
        self.track_gt_votes = defaultdict(Counter)
        self.track_gtid_votes = defaultdict(Counter)
        self.track_match_count = Counter()
        self.match_rows = []
        self.total_obs = 0
        self.total_iou_hits = 0

        # Tracking evaluation accumulators
        self.pred_to_gt_votes = defaultdict(Counter)
        self.gt_to_pred_votes = defaultdict(Counter)
        self.gt_frame_matches = defaultdict(list)
        self.association_rows = []

        # Frame-level OCR classification rows (filled after temporal voting)
        self.frame_ocr_rows = []

        if csv_path and os.path.exists(csv_path):
            self._load_csv(csv_path)
            self.available = True

    def _load_csv(self, csv_path):
        df = pd.read_csv(csv_path)
        df.columns = [str(c).strip().lower() for c in df.columns]
        # required columns: track_id, frame, jersey_number, x1, y1, x2, y2, ...
        ...
\end{lstlisting}
